% Provenance (never printed): sources/book2/BRIEFS.md, g11-05. Oligopoly and monopolistic competition, games,
% dominant strategies, Nash equilibrium, the prisoner's dilemma, cartels, repeated play, leniency. Real numbers:
% figdata/grade-11/05-oligopoly-and-games/ (EIA crude output 1979-1990; BLS real gasoline price in place of the
% unapproved World Bank crude price; Connor 2005 overcharges as a cited table). Ledger: 1980s oil glut, lysine,
% Connor 2005, EU leniency, the 2025 car-recycling decision. Kestria's quarries are invented.
\chapter{Oligopoly and a First Look at Games}\label{ch:g11:oligopoly-and-games}

Between 1981 and 1985, Saudi Arabia cut its output of crude oil from 9.81 million barrels a day to 3.39 million: by two-thirds. The other
members of the oil exporters' organisation, OPEC, which had agreed to share the cuts, produced 11.42 million barrels a day in 1981 and 11.92
million in 1985. Saudi Arabia was holding the price up almost alone. At the end of 1985 it gave up: in 1986 it pumped 4.87 million barrels a day,
and the price of oil fell from about 27 dollars a barrel to below 10.
% figdata: opec_points (sau 1981 9.81, 1985 3.39, 1986 4.87; others 1981 11.42, 1985 11.92); ledger: g11:opec-1980s

A market with a few large sellers is not like the markets of the last two chapters. Each seller's best choice depends on what the others do, and
each knows it. This chapter gives the tool for such situations, the theory of \omterm{def:g11:oligopoly-and-games:game}{games}, and uses it to explain why sellers would like to agree on high
prices, why their agreements tend to break, and what the law does about it.

\begin{openingchart}{Crude oil output of Saudi Arabia and of eleven other OPEC members, 1979--1990, in millions of barrels a day (top), and the price of a gallon of
regular petrol in the United States, in 2024 dollars (bottom).}{US Energy Information Administration, international energy statistics (crude oil including lease
condensate; the other members exclude Indonesia, with no data before 1984); US Bureau of Labor Statistics, average price data and CPI-U}
\begin{tikzpicture}
\begin{groupplot}[group style={group size=1 by 2, vertical sep=0.9cm}, omtime, width=0.84\linewidth, xmin=1979, xmax=1990,
  xtick={1979,1981,...,1989}]
\nextgroupplot[width=0.84\linewidth, height=5.0cm, ymin=0, ymax=22, ylabel={million barrels a day},
  legend style={at={(0.5,0.97)}, anchor=north, font=\scriptsize}]
\addplot[omThm, thick, mark=*, mark size=1pt] table[x=year, y=sau] {figdata/grade-11/05-oligopoly-and-games/out/opec.csv};
\addlegendentry{Saudi Arabia}
\addplot[omDef, thick, mark=*, mark size=1pt] table[x=year, y=others] {figdata/grade-11/05-oligopoly-and-games/out/opec.csv};
\addlegendentry{eleven other members}
\nextgroupplot[width=0.84\linewidth, height=3.6cm, ymin=2, ymax=5, ylabel={dollars a gallon}]
\addplot[black, thick, mark=*, mark size=1pt] table[x=year, y=gas] {figdata/grade-11/05-oligopoly-and-games/out/opec.csv};
\end{groupplot}
\end{tikzpicture}
\end{openingchart}

\begin{recall}
From the grades 6--9 volume: firms compete for buyers, and competition keeps prices down. From \cref{ch:g11:perfect-competition}: \omterm{def:g11:perfect-competition:economic-profit}{economic profit} and \omterm{def:g11:perfect-competition:free-entry}{free
entry}. From \cref{ch:g11:monopoly}: the \omterm{def:g11:monopoly:monopoly-output}{monopoly output}, where \omterm{def:g11:marginal-reasoning:marginal-benefit}{marginal revenue} equals \omterm{def:g11:marginal-reasoning:marginal-benefit}{marginal cost}, and the \omterm{def:g10:what-prices-do:dwl}{deadweight loss}.
\end{recall}

\section{Few sellers, or many different ones}\label{sec:g11:oligopoly-and-games:structures}

\begin{definition}[Oligopoly]\label{def:g11:oligopoly-and-games:oligopoly}
An \emph{oligopoly}\index{oligopoly} is a market with a few sellers, each large enough for its choices to change the price and the other sellers' profits.
\end{definition}

Mobile phone networks, aircraft makers, cement in a region and oil exporters are oligopolies. The opposite case has many sellers who each sell something a little
different.

\begin{definition}[Product differentiation and monopolistic competition]\label{def:g11:oligopoly-and-games:differentiation}
\emph{Product differentiation}\index{product differentiation} means that sellers' products differ in ways buyers care about: place, taste, brand, design.
\emph{Monopolistic competition}\index{monopolistic competition}\phantomsection\label{def:g11:oligopoly-and-games:monopolistic-competition} is a market with many
sellers of differentiated products and \omterm{def:g11:perfect-competition:free-entry}{free entry}.
\end{definition}

Restaurants, hairdressers and clothes shops fit. Each seller faces a demand curve of its own that slopes down, because some buyers prefer its product, so it
sets a price above \omterm{def:g11:marginal-reasoning:marginal-benefit}{marginal cost}, like a small monopoly. \omterm{def:g11:perfect-competition:free-entry}{Free entry} does the rest: while there are profits, new sellers open and take some of each seller's
buyers, until the demand curve of each just touches its average cost curve.

\begin{omfigure}
\begin{tikzpicture}[econ, x=0.55cm, y=0.14cm]
\draw[econ axes] (0,36) node[left, font=\small] {dalens} |- (13,0) node[below left=12pt and 0pt, font=\small] {$q$};
\draw[demand] (0,32) -- (10.67,0) node[above right, font=\scriptsize] {demand};
\draw[omMeth, thick] (0,32) -- (5.33,0) node[above, font=\scriptsize] {MR};
\draw[supply, domain=2.2:12.5, samples=50] plot (\x, {64/\x + \x}) node[right, font=\scriptsize] {ATC};
\draw[omThm, thick, domain=0:12] plot (\x, {2*\x}) node[right, font=\scriptsize] {MC};
\node[eqpoint] at (4,20) {}; \node[eqpoint] at (8,16) {};
\draw[guide] (0,20) -- (4,20) -- (4,0); \draw[guide] (8,16) -- (8,0); \draw[guide] (4,8) -- (0,8);
\foreach \y in {8,16,20,32} \draw (0,\y) -- (-0.15,\y) node[left, font=\footnotesize] {\y};
\node[below, font=\footnotesize] at (4,0) {4}; \node[below, font=\footnotesize] at (8,0) {8};
\end{tikzpicture}

{\small A seller in \omterm{def:g11:oligopoly-and-games:differentiation}{monopolistic competition} in the long run, with demand $p = 32 - 3q$, $\mathit{ATC} = 64/q + q$ and $\mathit{MC} = 2q$. It sells 4 units, where \omterm{def:g11:marginal-reasoning:marginal-benefit}{marginal
revenue} equals \omterm{def:g11:marginal-reasoning:marginal-benefit}{marginal cost} (8), at a price of 20, equal to its average cost: zero profit. The price is above \omterm{def:g11:marginal-reasoning:marginal-benefit}{marginal cost}, and the seller produces less than the 8
units at which its average cost would be lowest.}
\end{omfigure}

\section{Games}\label{sec:g11:oligopoly-and-games:games}

\begin{definition}[Game, strategy, payoff]\label{def:g11:oligopoly-and-games:game}
A \emph{game}\index{game} is a situation in which several players each choose an action, and each player's result depends on the choices of all. A
\emph{strategy}\index{strategy}\phantomsection\label{def:g11:oligopoly-and-games:strategy} is a choice open to a player; a player's
\emph{payoff}\index{payoff}\phantomsection\label{def:g11:oligopoly-and-games:payoff} is what she gets for each combination of strategies.
\end{definition}

With two players and two strategies each, the \omterm{def:g11:oligopoly-and-games:game}{game} fits in a table. Each cell gives the \omterm{def:g11:oligopoly-and-games:game}{payoffs} of the row player, then of the column player.

\begin{definition}[Dominant strategy]\label{def:g11:oligopoly-and-games:dominant-strategy}
A \omterm{def:g11:oligopoly-and-games:game}{strategy} is \emph{dominant}\index{dominant strategy} for a player if it gives her a higher \omterm{def:g11:oligopoly-and-games:game}{payoff} than any other \omterm{def:g11:oligopoly-and-games:game}{strategy}, whatever the other players do.
\end{definition}

\begin{definition}[Nash equilibrium]\label{def:g11:oligopoly-and-games:nash}
A \emph{Nash equilibrium}\index{Nash equilibrium} is a combination of strategies, one for each player, in which no player can raise her \omterm{def:g11:oligopoly-and-games:game}{payoff} by changing her
\omterm{def:g11:oligopoly-and-games:game}{strategy} while the others keep theirs.
\end{definition}

To find the Nash equilibria of a table, mark each player's best reply to each \omterm{def:g11:oligopoly-and-games:game}{strategy} of the other; a cell in which both \omterm{def:g11:oligopoly-and-games:game}{payoffs} are marked is an equilibrium.
If each player has a \omterm{def:g11:oligopoly-and-games:dominant-strategy}{dominant strategy}, the cell where both play it is the only equilibrium.

\section{The prisoner's dilemma and the cartel}\label{sec:g11:oligopoly-and-games:cartel}

\begin{definition}[Cartel and collusion]\label{def:g11:oligopoly-and-games:cartel}
A \emph{cartel}\index{cartel} is a group of firms that agree to limit their output or fix their prices, to earn together what a monopoly would.
\emph{Collusion}\index{collusion}\phantomsection\label{def:g11:oligopoly-and-games:collusion} is any such cooperation among rivals to raise prices, by an
explicit agreement or a tacit one.
\end{definition}

\begin{example}[Two quarries]\label{ex:g11:oligopoly-and-games:quarries}
Two quarries sell gravel; demand is $p = 120 - Q$ (thousand tonnes a year, kests a tonne) and each tonne costs 24 kests to produce. A monopoly would sell where $120 -
2Q = 24$, that is 48: so the \omterm{def:g11:oligopoly-and-games:cartel}{cartel} gives each quarry a quota of 24, the price is 72, and each earns $(72 - 24) \times 24 = 1152$ thousand kests. If one quarry
produces 32 while the other keeps to 24, the price is $120 - 56 = 64$: the first earns $40 \times 32 = 1280$ and the second $40 \times 24 = 960$. If both produce 32,
the price is 56 and each earns $32 \times 32 = 1024$.
\end{example}

\begin{omfigure}
\small
\begin{tabular}{l c c}
\toprule
& \multicolumn{2}{c}{quarry B} \\
quarry A & 24 (quota) & 32 \\
\midrule
24 (quota) & 1152, 1152 & 960, 1280 \\
32 & 1280, 960 & \textbf{1024, 1024} \\
\bottomrule
\end{tabular}

\medskip
{\small The two quarries' \omterm{def:g11:oligopoly-and-games:game}{game}: yearly profits in thousand kests (quarry A, quarry B). Producing 32 is a \omterm{def:g11:oligopoly-and-games:dominant-strategy}{dominant strategy} for each; the only \omterm{def:g11:oligopoly-and-games:nash}{Nash equilibrium}, in bold,
leaves both worse off than keeping to the quota.}
\end{omfigure}

\begin{definition}[Prisoner's dilemma]\label{def:g11:oligopoly-and-games:prisoners-dilemma}
A \emph{prisoner's dilemma}\index{prisoner's dilemma} is a \omterm{def:g11:oligopoly-and-games:game}{game} in which each player has a \omterm{def:g11:oligopoly-and-games:dominant-strategy}{dominant strategy}, but the equilibrium in which all play it gives every
player less than another combination would.
\end{definition}

A \omterm{def:g11:oligopoly-and-games:cartel}{cartel} is a \omterm{def:g11:oligopoly-and-games:prisoners-dilemma}{prisoner's dilemma}. Each member gains by selling more than its quota at the high price that the others' restraint keeps up, and if all reason so,
the \omterm{def:g11:oligopoly-and-games:cartel}{cartel} collapses. The OPEC members who produced above their quotas in the early 1980s, while Saudi Arabia cut, were playing the \omterm{def:g11:oligopoly-and-games:dominant-strategy}{dominant strategy}.

\begin{model}{The cartel as a prisoner's dilemma}
\assumes Each member of a \omterm{def:g11:oligopoly-and-games:cartel}{cartel} chooses its own output to maximise its own profit; the agreement cannot be enforced in court; the \omterm{def:g11:oligopoly-and-games:game}{game} is played once.
\predicts Each member produces more than its quota; the price falls towards the level without a \omterm{def:g11:oligopoly-and-games:cartel}{cartel}; \omterm{def:g11:oligopoly-and-games:cartel}{cartels} break down unless members can punish cheating.
\tested Consistent with it: OPEC in 1981--85, while Saudi Arabia cut its output by two-thirds to defend the price, the other members did not cut theirs; when Saudi Arabia stopped in
  1986, the price collapsed (opening chart).
\breaks When the same firms meet year after year and cheating is quickly seen, \omterm{def:g11:oligopoly-and-games:cartel}{cartels} can last: the lysine \omterm{def:g11:oligopoly-and-games:cartel}{cartel} of the 1990s met regularly and raised prices
  by 70\,\% within nine months. In laboratory experiments, players cooperate in a one-shot dilemma more often than the model predicts.
\end{model}
% ledger: g11:opec-1980s, g11:lysine; figdata: opec

\begin{record}[OPEC, 1981--86]
% ledger: g11:opec-1980s; figdata: opec_points, opec
In the early 1980s, the oil exporters' organisation set production quotas for its members to hold the price up as demand for oil fell. Several members cheated on
their quotas or refused them. Saudi Arabia, the largest producer, acted as the swing producer: it cut its own output to keep the total down, from 9.81 million
barrels a day in 1981 to 3.39 million in 1985. Then it abandoned that role and produced at full capacity. In 1986 the price of oil fell from about 27 dollars a barrel
to below 10, and the real price of petrol in the United States fell by about a quarter. Iran and Iraq were at war through most of the decade, which shaped their
output whatever the quotas said.
\end{record}
% figdata: opec (gas 1985 3.50, 1986 2.65)

\section{Playing again and again}\label{sec:g11:oligopoly-and-games:repeated}

\begin{definition}[Repeated game]\label{def:g11:oligopoly-and-games:repeated-game}
A \emph{repeated game}\index{repeated game} is the same \omterm{def:g11:oligopoly-and-games:game}{game} played by the same players several times, each player seeing what the others did before choosing again.
\end{definition}

In a \omterm{def:g11:oligopoly-and-games:repeated-game}{repeated game}, a player can make her future choices depend on the others' past ones. A \emph{trigger \omterm{def:g11:oligopoly-and-games:game}{strategy}} keeps to the quota as long as the other does, and
produces the high output for ever after the first time the other cheats. Against it, cheating brings a gain once and a loss every year after. In the quarries' \omterm{def:g11:oligopoly-and-games:game}{game},
cheating with the best possible output brings 144 thousand kests in the year of cheating (the weekend problem computes it) and costs $1152 - 1024 = 128$ in each later
year: two years of punishment remove the gain. The more the players care about the future, and the faster cheating is seen, the easier it is to sustain a \omterm{def:g11:oligopoly-and-games:cartel}{cartel}.

\section{Cartels and the law}\label{sec:g11:oligopoly-and-games:law}

\omterm{def:g11:oligopoly-and-games:cartel}{Cartels} are illegal in most countries, and their members are fined, and in some countries jailed. In the lysine case of the 1990s, an executive of one member firm
informed the FBI, the American federal police; the firm was fined 100 million dollars, and three of its executives went to prison.
% ledger: g11:lysine

\begin{definition}[Leniency programme]\label{def:g11:oligopoly-and-games:leniency}
A \emph{leniency programme}\index{leniency programme} cancels or reduces the fine of a \omterm{def:g11:oligopoly-and-games:cartel}{cartel} member that reports the \omterm{def:g11:oligopoly-and-games:cartel}{cartel} to the competition authority and helps
prove it.
\end{definition}

The European Union has run one since 1996; under its notice of 2006, the first member to report gets full immunity, and later members who bring evidence of real
value get reductions of 30 to 50\,\%, 20 to 30\,\%, and up to 20\,\%. Leniency starts a race to report: each member fears that another will report first, and once one is
expected to report, every member's best reply is to report too.
% ledger: g11:eu-leniency

\begin{dated}{2026}
% ledger: g11:eu-car-recycling-2025
On 1 April 2025 the European Commission announced fines of 458 million euros on 15 carmakers and their trade association, which it found had coordinated for more than
15 years to avoid paying for the recycling of old cars and not to compete on how recyclable their cars were. The carmaker that revealed the \omterm{def:g11:oligopoly-and-games:cartel}{cartel} avoided a fine of
about 35 million euros; another had its fine halved, to 75 million euros, for cooperating.
\end{dated}

How much do \omterm{def:g11:oligopoly-and-games:cartel}{cartels} raise prices? A survey of hundreds of studies, with 699 estimates, found a median overcharge, the rise of the price above what it would have been
without the \omterm{def:g11:oligopoly-and-games:cartel}{cartel}, of 25\,\%, and a mean of 42\,\%, pulled up by a few very large ones.
% ledger: g11:connor2005-all

\begin{omchart}{Median long-run cartel overcharge, in percent of the price that would have prevailed without the cartel, by kind of cartel.}{Connor (2005), survey of
published overcharge estimates}
\begin{tikzpicture}
\begin{axis}[omtime, ybar, width=0.62\linewidth, height=5.0cm, ymin=0, ymax=40, bar width=18pt, symbolic x coords={all cartels,domestic,international,successful},
  enlarge x limits=0.18, xtick=data, ylabel={percent}, nodes near coords, nodes near coords style={font=\scriptsize}, xmajorgrids=false, x tick label style={/tikz/font=\scriptsize}]
\addplot[fill=omDef!60, draw=omDef] table[x=group, y=median] {figdata/grade-11/05-oligopoly-and-games/out/overcharge.csv};
\end{axis}
\end{tikzpicture}
\end{omchart}
% ledger: g11:connor2005-all, g11:connor2005-domestic, g11:connor2005-international, g11:connor2005-successful

Not every \omterm{def:g11:oligopoly-and-games:game}{game} is a dilemma. In a \emph{coordination \omterm{def:g11:oligopoly-and-games:game}{game}} the players want to make the same choice: drive on the same side of the road, use the same plug. There
are then several Nash equilibria, and the theory alone does not say which one is reached; custom, a law or a first mover does.

\section{Exercises}

\begin{exercise}[$\star$]\label{exo:g11:oligopoly-and-games:1}
In this \omterm{def:g11:oligopoly-and-games:game}{game}, \omterm{def:g11:oligopoly-and-games:game}{payoffs} are (row, column). Row: Top gives (4, 1) against Left and (2, 3) against Right; Bottom gives (3, 0) and (1, 2). Does either player have a \omterm{def:g11:oligopoly-and-games:dominant-strategy}{dominant
strategy}? What is the outcome?
\end{exercise}

\begin{exercise}[$\star$]\label{exo:g11:oligopoly-and-games:2}
Two cafés each choose a high or a low price. Profits (A, B): both high (10, 10); A high and B low (4, 12); A low and B high (12, 4); both low (6, 6). Find each café's
\omterm{def:g11:oligopoly-and-games:dominant-strategy}{dominant strategy} and the outcome.
\end{exercise}

\begin{exercise}[$\star$]\label{exo:g11:oligopoly-and-games:3}
\omterm{def:g11:oligopoly-and-games:game}{Payoffs} (row, column): (Top, Left) (5, 2); (Top, Right) (3, 1); (Bottom, Left) (4, 1); (Bottom, Right) (2, 3). Which player has a \omterm{def:g11:oligopoly-and-games:dominant-strategy}{dominant strategy}? Find the \omterm{def:g11:oligopoly-and-games:nash}{Nash
equilibrium}.
\end{exercise}

\begin{exercise}[$\star$]\label{exo:g11:oligopoly-and-games:4}
Two firms choose a standard, A or B. \omterm{def:g11:oligopoly-and-games:game}{Payoffs}: both A (3, 3); both B (1, 1); different standards (0, 0). Find all the Nash equilibria.
\end{exercise}

\begin{exercise}[$\star$]\label{exo:g11:oligopoly-and-games:5}
Read the opening chart. By how much did Saudi output fall between 1981 and 1985, and the others' output change? What happened to both and to the price in 1986?
\end{exercise}

\begin{exercise}[$\star\star$]\label{exo:g11:oligopoly-and-games:6}
Demand is $p = 100 - Q$ and two firms each have a \omterm{def:g11:marginal-reasoning:marginal-benefit}{marginal cost} of 10. Each can produce 22.5 (half the \omterm{def:g11:monopoly:monopoly-output}{monopoly output}) or 30. Build the \omterm{def:g11:oligopoly-and-games:game}{payoff} table and find the
equilibrium.
\end{exercise}

\begin{exercise}[$\star\star$]\label{exo:g11:oligopoly-and-games:7}
Same question with $p = 130 - Q$, a \omterm{def:g11:marginal-reasoning:marginal-benefit}{marginal cost} of 10, and outputs of 30 or 40.
\end{exercise}

\begin{exercise}[$\star\star$]\label{exo:g11:oligopoly-and-games:8}
In \cref{exo:g11:oligopoly-and-games:7}, the other firm keeps to 30. Find, with a derivative, the output that maximises the first firm's profit, and its gain over
keeping to 30.
\end{exercise}

\begin{exercise}[$\star\star$]\label{exo:g11:oligopoly-and-games:9}
Two \omterm{def:g11:oligopoly-and-games:cartel}{cartel} members each choose to stay silent or to report the \omterm{def:g11:oligopoly-and-games:cartel}{cartel}. Silent, each expects a profit of 60 (including the risk of being caught). If one reports, it gets
immunity and earns 0, and the other is fined: $-200$. If both report, each gets half the fine: $-100$. Find the Nash equilibria. How does the \omterm{def:g11:oligopoly-and-games:leniency}{leniency programme} threaten
the \omterm{def:g11:oligopoly-and-games:cartel}{cartel}?
\end{exercise}

\begin{exercise}[$\star\star$]\label{exo:g11:oligopoly-and-games:10}
Read the overcharge chart. A \omterm{def:g11:oligopoly-and-games:cartel}{cartel} with the median overcharge of an international \omterm{def:g11:oligopoly-and-games:cartel}{cartel} sells goods worth 660 million euros at its prices. How much did buyers pay
because of the \omterm{def:g11:oligopoly-and-games:cartel}{cartel}?
\end{exercise}

\begin{exercise}[$\star\star$]\label{exo:g11:oligopoly-and-games:11}
Find all the Nash equilibria: (Top, Left) (3, 2); (Top, Right) (0, 0); (Bottom, Left) (0, 0); (Bottom, Right) (2, 3). Which would each player prefer?
\end{exercise}

\begin{exercise}[$\star\star\star$]\label{exo:g11:oligopoly-and-games:12}
What evidence contradicts the one-shot model of the \omterm{def:g11:oligopoly-and-games:cartel}{cartel}? Use the lysine \omterm{def:g11:oligopoly-and-games:cartel}{cartel} and laboratory experiments, and say which assumption of the model fails in each.
\end{exercise}

\begin{exercise}[$\star\star\star$]\label{exo:g11:oligopoly-and-games:13}
In a \omterm{def:g11:oligopoly-and-games:repeated-game}{repeated game}, cooperation brings $\pi_c$ each year, cheating once brings $\pi_d$ then $\pi_n$ for ever, and each player values a kest next year at $\delta$ kests now
($0 < \delta < 1$). Admitting that $\delta + \delta^2 + \dots = \delta/(1 - \delta)$, show that a trigger \omterm{def:g11:oligopoly-and-games:game}{strategy} deters cheating if $\pi_d - \pi_c \leq (\pi_c - \pi_n)\,\delta/(1 -
\delta)$. Find the smallest $\delta$ for the two quarries.
\end{exercise}

\begin{exercise}[$\star\star\star$]\label{exo:g11:oligopoly-and-games:14}
Explain why a firm in \omterm{def:g11:oligopoly-and-games:differentiation}{monopolistic competition} makes zero profit in the long run although its price is above its \omterm{def:g11:marginal-reasoning:marginal-benefit}{marginal cost}, and why it produces less than the output at
which its average cost is lowest.
\end{exercise}

\begin{exercise}[$\star\star\star$]\label{exo:g11:oligopoly-and-games:15}
Which side of the road to drive on is a coordination \omterm{def:g11:oligopoly-and-games:game}{game}. Write its \omterm{def:g11:oligopoly-and-games:game}{payoff} table for two drivers, find its equilibria, and explain why a law settles it while no law could
settle a \omterm{def:g11:oligopoly-and-games:prisoners-dilemma}{prisoner's dilemma} without a penalty.
\end{exercise}

\section{Problem: Two Quarries in Kestria}

\begin{problem}[{Weekend problem --- a cartel, cheating and punishment}]
\label{pb:g11:oligopoly-and-games:1}
Two quarries supply all the gravel in a Kestrian valley. Demand is $p = 120 - Q$, with $Q$ in thousand tonnes a year and $p$ in kests a tonne; each quarry's cost is 24
kests a tonne, with no fixed cost. Profits are in thousand kests a year.

\medskip
\noindent\textbf{Part I --- The cartel.}
\begin{enumerate}
\item Find the output and price that a single owner of both quarries would choose.
\item The \omterm{def:g11:oligopoly-and-games:cartel}{cartel} splits that output equally. Give each quarry's quota and profit.
\item Compute the \omterm{def:g10:what-prices-do:dwl}{deadweight loss} of the \omterm{def:g11:oligopoly-and-games:cartel}{cartel}, compared with a price equal to \omterm{def:g11:marginal-reasoning:marginal-benefit}{marginal cost}.
\item Quarry B keeps to its quota of 24. Write quarry A's profit as a function of its own output $q_A$.
\item Find, with the derivative, the output that maximises it, the price and A's profit.
\end{enumerate}

\medskip
\noindent\textbf{Part II --- Without a cartel.}
\begin{enumerate}[resume]
\item For any output $q_B$ of quarry B, find A's best output (its best reply) as a function of $q_B$.
\item Find the outputs at which each quarry's output is the best reply to the other's.
\item Compute the price and each quarry's profit there.
\item Build the \omterm{def:g11:oligopoly-and-games:game}{payoff} table with the strategies 24 and 32 for each quarry.
\item Show that the \omterm{def:g11:oligopoly-and-games:game}{game} is a \omterm{def:g11:oligopoly-and-games:prisoners-dilemma}{prisoner's dilemma}.
\end{enumerate}

\medskip
\noindent\textbf{Part III --- Cheating and punishment.}
\begin{enumerate}[resume]
\item Compute A's gain from its best cheating output (question 5) in the year it cheats.
\item What does B earn that year?
\item After cheating is seen, both produce 32 for ever. How much does A lose in each later year, compared with the \omterm{def:g11:oligopoly-and-games:cartel}{cartel}?
\item Without discounting, how many years of punishment remove the gain?
\item With a kest next year worth $1/1.1$ kests now, are two years of punishment enough? Compute the discounted loss.
\end{enumerate}

\medskip
\noindent\textbf{Part IV --- A ten-year horizon.}
\begin{enumerate}[resume]
\item The quarries' licences end in ten years. In the tenth year, is there any reason to keep to the quota? Explain.
\item What does this imply for the ninth year, and for the years before?
\item In practice, firms rarely know when their \omterm{def:g11:oligopoly-and-games:game}{game} ends. Why does that help the \omterm{def:g11:oligopoly-and-games:cartel}{cartel}?
\item Under a \omterm{def:g11:oligopoly-and-games:leniency}{leniency programme}, a quarry that reports the \omterm{def:g11:oligopoly-and-games:cartel}{cartel} pays no fine. How does that change the comparison of question 14?
\item State the result of the problem: the gain from cheating on the \omterm{def:g11:oligopoly-and-games:cartel}{cartel}, and the punishment that removes it.
\end{enumerate}
\end{problem}
